% Generated draft: requires mathematical and source review.
\documentclass[11pt]{article}
\usepackage[margin=25mm]{geometry}
\usepackage{fontspec}
\setmainfont{Latin Modern Roman}
\usepackage{amsmath,amssymb,mathtools,unicode-math}
\setmathfont{Latin Modern Math}
\usepackage{enumitem,needspace,xurl,hyperref}
\setlength{\parindent}{0pt}
\setlength{\parskip}{5pt}
\setlength{\emergencystretch}{3em}
\newcommand{\N}{\mathbb{N}}
\newcommand{\Z}{\mathbb{Z}}
\newcommand{\Q}{\mathbb{Q}}
\newcommand{\R}{\mathbb{R}}
\newcommand{\C}{\mathbb{C}}
\newcommand{\F}{\mathbb{F}}
\newcommand{\A}{\mathbb{A}}
\newcommand{\PP}{\mathbb{P}}
\newcommand{\E}{\mathbb{E}}
\newcommand{\eps}{\varepsilon}
\newcommand{\dd}{\,\mathrm{d}}
\begin{document}
\section*{UnsolvedMath ID: 700482}
\textbf{Source problem code:} catalogue entry
\subsection{Definitions and mathematical statement}
A group \(G\) is soluble if its derived series, defined by \(G^{(0)}=G\) and \(G^{(n+1)}=[G^{(n)},G^{(n)}]\), reaches the trivial group after finitely many steps. Let \(\mathfrak S\) denote the class of all soluble groups in the ambient universe intended by the source. The supplied record does not explicitly specify whether all groups are required to be finite; this domain convention remains unresolved.

A Fitting class is a nonempty class of groups, invariant under isomorphism, closed under taking normal subgroups and under normal products: if \(N_1,N_2\trianglelefteq G\) belong to the class, then \(N_1N_2\) belongs to it, within the intended ambient universe. For a Fitting class \(\mathfrak F\), write \(\mathfrak F_*\) for the bottom member of the Lockett section determined by \(\mathfrak F\). This retains the source's named construction; the supplied record gives neither a definition of a Lockett section nor an independent construction of \(\mathfrak F_*\). The phrase ``bottom group'' in the source is interpreted here as a bottom class, since the displayed equalities are equalities of classes.

The Lockett condition for \(\mathfrak F\) is
\[
\mathfrak F\cap\mathfrak S_*=\mathfrak F_*.
\]
The question is whether, for every pair of Fitting classes \(\mathfrak X,\mathfrak Y\subseteq\mathfrak S\),
\[
\left(\mathfrak X\cap\mathfrak S_*=\mathfrak X_*\ \text{ and }\ \mathfrak Y\cap\mathfrak S_*=\mathfrak Y_*\right)
\quad\Longrightarrow\quad
(\mathfrak X\cap\mathfrak Y)\cap\mathfrak S_*=(\mathfrak X\cap\mathfrak Y)_*.
\]
Thus the intended assertion is closure under binary intersection of the Fitting classes satisfying the Lockett condition, subject to the unresolved ambient convention and the source's Lockett-section terminology.

The supplied editors' comment, dated 2001, reports an affirmative answer when both \(\mathfrak X\) and \(\mathfrak Y\) are local Fitting classes, citing A. Grytczuk and N. T. Vorob'ev (1994). This is a source-reported special case; the record does not define locality.
\subsection{Short English statement}
Must the intersection of any two Fitting classes of soluble groups satisfying the Lockett condition also satisfy that condition? The supplied editors' comment reports an affirmative answer when both classes are local. The finiteness convention and definitions of Lockett sections and locality require confirmation.
\subsection{Sources}
\begin{itemize}
\item Local UnsolvedMath registry, permanent ID \(700482\), exact source problem code \texttt{catalogue entry}, title \emph{Kourovka Notebook Problem 8.30}. The supplied registry labels the problem \emph{open}; this label is not an independently verified mathematical or literature assessment.
\item H. Lausch, Problem 8.30, \emph{The Kourovka Notebook: Unsolved Problems in Group Theory}, 8th issue (1982). The supplied background identifies the 21st-issue edition (2026), arXiv:1401.0300v48. The collection description names E. I. Khukhro and V. D. Mazurov as editors.
\item Editors' comment (2001): the answer is affirmative for local Fitting classes, attributed to A. Grytczuk and N. T. Vorob'ev, \emph{Tsukuba J. Math.} \textbf{18}, no. 1 (1994), 63--67. The supplied record gives no article title or link.
\item Source provenance: supplied upstream revision \texttt{d371ee917ac9647fbb06cd6fc529deec6f3dda82}; registry update timestamp \texttt{2026-10-16T14:00:00Z}. This timestamp is registry metadata, not a stated literature-verification date. No independent source consultation was performed.
\end{itemize}
\end{document}
